\begin{table}[!htbp]
\centering
\footnotesize
\def\sym#1{\ifmmode^{#1}\else\(^{#1}\)\fi}
\caption{Politician addresses appearing in stealer-log corpora}
\label{tab:stealer_logs}
\begin{adjustbox}{max width=\linewidth}
\begin{tabular}{llrrl}
\toprule\toprule
Corpus & Breach date & Addresses & \% of sample & 95\% CI \\
\midrule
TelegramCombolists & 2024-05 & 163 & 1.32\% & [1.13, 1.53] \\
TrikSpamBotnet & 2018-06 &  95 & 0.77\% & [0.63, 0.94] \\
NazApi & 2023-09 &  63 & 0.51\% & [0.40, 0.65] \\
StealerLogsJan2025 & 2025-01 &  40 & 0.32\% & [0.24, 0.44] \\
TelegramStealerLogs & 2024-07 &  30 & 0.24\% & [0.17, 0.35] \\
RedLineStealer & 2021-12 &   2 & 0.02\% & [0.00, 0.06] \\
Emotet & 2021-01 &   0 & 0.00\% & [0.00, 0.03] \\
NamelessMalware & 2020-01 &   0 & 0.00\% & [0.00, 0.03] \\
Qakbot & 2023-08 &   0 & 0.00\% & [0.00, 0.03] \\
\textit{Any stealer corpus} &  & 286 & 2.31\% & [2.06, 2.59] \\
\quad personal addresses &  & 166 & 5.51\% & [4.75, 6.38] \\
\quad official addresses &  & 120 & 1.28\% & [1.07, 1.53] \\
\bottomrule
\end{tabular}
\end{adjustbox}
\caption*{\scriptsize \emph{Note:} HIBP loads some infostealer and malware corpora into the ordinary breach dataset; several originate in documented law-enforcement seizures (Emotet, from the FBI and Dutch NHTCU takedown; Qakbot, from the US Department of Justice operation). Unlike the rest of the catalogue, appearing in one indicates credentials captured from an infected machine rather than a third party holding the address. Intervals are Wilson binomial intervals. Per-corpus counts sum to more than the total because addresses appear in more than one corpus. \textbf{These are lower bounds even by the standards of the rest of this paper}: HIBP's dedicated stealer-log endpoint indexes more data, and it cannot be queried for this sample at any subscription tier, because access requires verified control of the address or of its email domain.}
\end{table}
